\documentclass[10pt,a4paper]{amsart}
\usepackage{amsmath,amssymb,amsthm,graphicx}
% Retained minimal preamble for the Stacks Project chapter more-algebra.tex (More on Algebra).
% The chapter file pinned-original-source.tex carries only \input{preamble} on its first line, so no
% preamble text lives inside the chapter itself.  The Stacks Project uses one single shared preamble for
% every chapter; the snapshot stages it identically as sources/stacks-chow/preamble.tex and
% sources/stacks-etale-cohomology/preamble.tex.  The block below is copied mechanically, line for line,
% from sources/stacks-chow/preamble.tex lines 14-18 (the Xy-pic setup) and 154-217 (the native theorem
% declarations together with the mathematical macros).  Nothing was added, removed or redefined; the
% wrapper class and the amsmath/amssymb/amsthm/graphicx base packages are supplied by the assembler.
% Chapter: more-algebra.tex -- More on Algebra.
% For commutative diagrams we use Xy-pic
\usepackage[all]{xy}

% We use 2cell for 2-commutative diagrams.
\xyoption{2cell}

% Theorem environments.
%
\theoremstyle{plain}
\newtheorem{theorem}[subsection]{Theorem}
\newtheorem{proposition}[subsection]{Proposition}
\newtheorem{lemma}[subsection]{Lemma}

\theoremstyle{definition}
\newtheorem{definition}[subsection]{Definition}
\newtheorem{example}[subsection]{Example}
\newtheorem{exercise}[subsection]{Exercise}
\newtheorem{situation}[subsection]{Situation}

\theoremstyle{remark}
\newtheorem{remark}[subsection]{Remark}
\newtheorem{remarks}[subsection]{Remarks}

\numberwithin{equation}{subsection}

% Macros
%
\def\lim{\mathop{\mathrm{lim}}\nolimits}
\def\colim{\mathop{\mathrm{colim}}\nolimits}
\def\Spec{\mathop{\mathrm{Spec}}}
\def\Hom{\mathop{\mathrm{Hom}}\nolimits}
\def\Ext{\mathop{\mathrm{Ext}}\nolimits}
\def\SheafHom{\mathop{\mathcal{H}\!\mathit{om}}\nolimits}
\def\SheafExt{\mathop{\mathcal{E}\!\mathit{xt}}\nolimits}
\def\Sch{\mathit{Sch}}
\def\Mor{\mathop{\mathrm{Mor}}\nolimits}
\def\Ob{\mathop{\mathrm{Ob}}\nolimits}
\def\Sh{\mathop{\mathit{Sh}}\nolimits}
\def\NL{\mathop{N\!L}\nolimits}
\def\CH{\mathop{\mathrm{CH}}\nolimits}
\def\proetale{{pro\text{-}\acute{e}tale}}
\def\etale{{\acute{e}tale}}
\def\QCoh{\mathit{QCoh}}
\def\Ker{\mathop{\mathrm{Ker}}}
\def\Im{\mathop{\mathrm{Im}}}
\def\Coker{\mathop{\mathrm{Coker}}}
\def\Coim{\mathop{\mathrm{Coim}}}

% Boxtimes
%
\DeclareMathSymbol{\boxtimes}{\mathbin}{AMSa}{"02}

%
% Macros for moduli stacks/spaces
%
\def\QCohstack{\mathcal{QC}\!\mathit{oh}}
\def\Cohstack{\mathcal{C}\!\mathit{oh}}
\def\Spacesstack{\mathcal{S}\!\mathit{paces}}
\def\Quotfunctor{\mathrm{Quot}}
\def\Hilbfunctor{\mathrm{Hilb}}
\def\Curvesstack{\mathcal{C}\!\mathit{urves}}
\def\Polarizedstack{\mathcal{P}\!\mathit{olarized}}
\def\Complexesstack{\mathcal{C}\!\mathit{omplexes}}
% \Pic is the operator that assigns to X its picard group, usage \Pic(X)
% \Picardstack_{X/B} denotes the Picard stack of X over B
% \Picardfunctor_{X/B} denotes the Picard functor of X over B
\def\Pic{\mathop{\mathrm{Pic}}\nolimits}
\def\Picardstack{\mathcal{P}\!\mathit{ic}}
\def\Picardfunctor{\mathrm{Pic}}
\def\Deformationcategory{\mathcal{D}\!\mathit{ef}}

% --- appended 2026-10-12 (Chief): the project own definitions of the reference/slogan/history
% display environments, copied mechanically from the shared Stacks preamble lines 8-12
% (\usepackage{verbatim} plus the three \newenvironment lines, which the chapter preamble keeps
% verbatim from the project preamble). Nothing added, removed or redefined: these environments
% carry no mathematical content and the project itself typesets them as comments. Without these
% three lines a counted statement whose own text uses them fails with Environment slogan undefined.
\usepackage{verbatim}
\newenvironment{reference}{\comment}{\endcomment}
\newenvironment{slogan}{\comment}{\endcomment}
\newenvironment{history}{\comment}{\endcomment}
\usepackage{cleveref}
\crefname{originallemma}{original lemma}{original lemmas}
\Crefname{originallemma}{Original lemma}{Original lemmas}
\makeatletter
\newlabel{more-algebra-proposition-ubiquity-J-2}{{07PJ}{source}{}{sourcepointer-1}{}}
\newlabel{more-algebra-proposition-ubiquity-J-2@cref}{{[originallemma][0][]07PJ}{[1][0][]source}{}{}{}}
\makeatother
\crefname{originallemma}{original lemma}{original lemmas}
\Crefname{originallemma}{Original lemma}{Original lemmas}
\makeatletter
\newlabel{lemma-J-2}{{07PC}{source}{}{sourcepointer-2}{}}
\newlabel{lemma-J-2@cref}{{[originallemma][0][]07PC}{[1][0][]source}{}{}{}}
\makeatother
\crefname{originallemma}{original lemma}{original lemmas}
\Crefname{originallemma}{Original lemma}{Original lemmas}
\makeatletter
\newlabel{algebra-lemma-integral-dim-up}{{00OJ}{source}{}{sourcepointer-3}{}}
\newlabel{algebra-lemma-integral-dim-up@cref}{{[originallemma][0][]00OJ}{[1][0][]source}{}{}{}}
\makeatother
\crefname{originallemma}{original lemma}{original lemmas}
\Crefname{originallemma}{Original lemma}{Original lemmas}
\makeatletter
\newlabel{algebra-lemma-criterion-normal}{{031S}{source}{}{sourcepointer-4}{}}
\newlabel{algebra-lemma-criterion-normal@cref}{{[originallemma][0][]031S}{[1][0][]source}{}{}{}}
\makeatother
\crefname{originallemma}{original lemma}{original lemmas}
\Crefname{originallemma}{Original lemma}{Original lemmas}
\makeatletter
\newlabel{algebra-lemma-quotient-complete-local}{{0325}{source}{}{sourcepointer-5}{}}
\newlabel{algebra-lemma-quotient-complete-local@cref}{{[originallemma][0][]0325}{[1][0][]source}{}{}{}}
\makeatother
\crefname{originallemma}{original lemma}{original lemmas}
\Crefname{originallemma}{Original lemma}{Original lemmas}
\makeatletter
\newlabel{lemma-intersection-regular-with-closed}{{07P9}{source}{}{sourcepointer-6}{}}
\newlabel{lemma-intersection-regular-with-closed@cref}{{[originallemma][0][]07P9}{[1][0][]source}{}{}{}}
\makeatother
\crefname{originallemma}{original lemma}{original lemmas}
\Crefname{originallemma}{Original lemma}{Original lemmas}
\makeatletter
\newlabel{lemma-complete-Noetherian-domain-J-0}{{07PI}{source}{}{sourcepointer-7}{}}
\newlabel{lemma-complete-Noetherian-domain-J-0@cref}{{[originallemma][0][]07PI}{[1][0][]source}{}{}{}}
\makeatother
\title{Fields, complete local rings, the integers, one-dimensional local and Nagata rings, and Dedekind domains of characteristic zero are J-2}
\author{The Stacks Project Authors}
\date{stacks-project:more-algebra.tex:local-snapshot-20261012; GNU FDL 1.2 or later}
\begin{document}
\maketitle
\setcounter{subsection}{1}
\noindent\textit{Editorial note.} All delivered mathematical text is the author's own, copied from the pinned Stacks Project chapter file more-algebra.tex; no statement or proof step was rewritten, reordered or removed.\par\medskip
\noindent\textit{Editorial note.} The chapter file pinned-original-source.tex carries only \textbackslash{}input\{preamble\} on its first line, so no preamble text lives inside the chapter. The retained file sources/stacks-more-algebra/retained-preamble.tex is a mechanical, line-for-line copy of the project's single shared preamble as staged in the snapshot (Xy-pic setup plus theorem declarations and mathematical macros). Nothing was added, removed or redefined, and the excerpt serves only as preamble, never as a span coordinate.\par\medskip
\noindent\textit{Editorial note.} Both counted spans start at the \textbackslash{}begin line of their environment, so the statement and the author's proof are delivered whole; surrounding prose of the section is deliberately outside the counted spans and is carried as standard prerequisites instead.\par\medskip
\noindent\textit{Editorial note.} The citation fields (source\_id, source\_version, authors, source\_url, license.name, source\_title) are identical to those of every other card this lane has produced from this chapter, so the batch selector treats them as one source.\par\medskip
\noindent\textit{Allowed standard prerequisites (editorial).}\begin{itemize}
\item The two criteria for a ring to be J-2: the condition on finite purely inseparable extensions of residue fields, and the condition on finite ring extensions being J-1. Lemma \textbackslash{}ref\{lemma-J-2\}, cited at proof lines 11959-11960, 11978-11979, and 11987-11988. Provide the criterion that is verified for each of the seven classes.
\item The integral closure of a one-dimensional Nagata domain in a finite extension is a normal domain of dimension one, hence regular. Algebra, Lemma \textbackslash{}ref\{algebra-lemma-integral-dim-up\} and Algebra, Lemma \textbackslash{}ref\{algebra-lemma-criterion-normal\}, cited at proof lines 11968-11971. Settle the case of a minimal prime in a one-dimensional Nagata ring.
\item A finite extension of a complete local ring is a product of complete local rings. Algebra, Lemma \textbackslash{}ref\{algebra-lemma-quotient-complete-local\}, cited at proof lines 11980-11981. Reduces the complete local case to the case of a domain via the regular-locus criterion.
\item A Noetherian complete local domain is J-0, and regularity is preserved when the regular locus meets every closed subset in a nonempty open subset. Lemma \textbackslash{}ref\{lemma-complete-Noetherian-domain-J-0\} and Lemma \textbackslash{}ref\{lemma-intersection-regular-with-closed\}, cited at proof lines 11982-11984. Finish the complete local case.
\end{itemize}\medskip
\ifdefined\hypertarget\hypertarget{sourcepointer-1}{}\fi
\noindent\textit{Source reference.} Original lemma 07PJ: focused-seeds/AGP/sources/stacks-snapshot-20261011/tags.txt, source lines 9799--9799. The lemma's own canonical tag, quoted so that readers can resolve the label inside the delivered text; it points at the counted statement itself and carries no proof content.\par\medskip
\ifdefined\hypertarget\hypertarget{sourcepointer-2}{}\fi
\noindent\textit{Source reference.} Original lemma 07PC: focused-seeds/AGP/sources/stacks-snapshot-20261011/tags.txt, source lines 9792--9792. Navigation pointer to a statement the counted proof refers to but does not contain; nothing load-bearing is delegated to it.\par\medskip
\ifdefined\hypertarget\hypertarget{sourcepointer-3}{}\fi
\noindent\textit{Source reference.} Original lemma 00OJ: focused-seeds/AGP/sources/stacks-snapshot-20261011/tags.txt, source lines 897--897. Navigation pointer to a statement the counted proof refers to but does not contain; nothing load-bearing is delegated to it.\par\medskip
\ifdefined\hypertarget\hypertarget{sourcepointer-4}{}\fi
\noindent\textit{Source reference.} Original lemma 031S: focused-seeds/AGP/sources/stacks-snapshot-20261011/tags.txt, source lines 3995--3995. Navigation pointer to a statement the counted proof refers to but does not contain; nothing load-bearing is delegated to it.\par\medskip
\ifdefined\hypertarget\hypertarget{sourcepointer-5}{}\fi
\noindent\textit{Source reference.} Original lemma 0325: focused-seeds/AGP/sources/stacks-snapshot-20261011/tags.txt, source lines 4008--4008. Navigation pointer to a statement the counted proof refers to but does not contain; nothing load-bearing is delegated to it.\par\medskip
\ifdefined\hypertarget\hypertarget{sourcepointer-6}{}\fi
\noindent\textit{Source reference.} Original lemma 07P9: focused-seeds/AGP/sources/stacks-snapshot-20261011/tags.txt, source lines 9789--9789. Navigation pointer to a statement the counted proof refers to but does not contain; nothing load-bearing is delegated to it.\par\medskip
\ifdefined\hypertarget\hypertarget{sourcepointer-7}{}\fi
\noindent\textit{Source reference.} Original lemma 07PI: focused-seeds/AGP/sources/stacks-snapshot-20261011/tags.txt, source lines 9798--9798. Navigation pointer to a statement the counted proof refers to but does not contain; nothing load-bearing is delegated to it.\par\medskip
\begin{proposition}
\label{proposition-ubiquity-J-2}
The following types of rings are J-2:
\begin{enumerate}
\item fields,
\item Noetherian complete local rings,
\item $\mathbf{Z}$,
\item Noetherian local rings of dimension $1$,
\item Nagata rings of dimension $1$,
\item Dedekind domains with fraction field of characteristic zero,
\item finite type ring extensions of any of the above.
\end{enumerate}
\end{proposition}

\begin{proof}
For cases (1), (3), (5), and (6) this is proved by checking
condition (4) of Lemma \ref{lemma-J-2}. We will only do this
in case $R$ is a Nagata ring of dimension $1$. Let $\mathfrak p \subset R$
be a prime ideal and let $L/\kappa(\mathfrak p)$ be a finite
purely inseparable extension. If $\mathfrak p \subset R$ is a
maximal ideal, then $R \to L$ is finite and $L$ is a regular ring
and we've checked the condition. If $\mathfrak p \subset R$ is a
minimal prime, then the Nagata condition insures that the
integral closure $R' \subset L$ of $R$ in $L$ is finite over $R$.
Then $R'$ is a normal domain of dimension $1$
(Algebra, Lemma \ref{algebra-lemma-integral-dim-up})
hence regular (Algebra, Lemma \ref{algebra-lemma-criterion-normal})
and we've checked the condition in this case as well.

\medskip\noindent
For case (2), we will use condition (3) of Lemma \ref{lemma-J-2}.
Let $R$ be a Noetherian complete local ring.
Note that if $R \to R'$ is finite, then $R'$ is a product of
Noetherian complete local rings, see
Algebra, Lemma \ref{algebra-lemma-quotient-complete-local}.
Thus by Lemma \ref{lemma-intersection-regular-with-closed} it
suffices to prove that a Noetherian complete local ring
which is a domain is J-0, which is
Lemma \ref{lemma-complete-Noetherian-domain-J-0}.

\medskip\noindent
For case (4), we also use condition (3) of Lemma \ref{lemma-J-2}.
Namely, if $R$ is a local Noetherian ring of dimension $1$ and
$R \to R'$ is finite, then $\Spec(R')$ is finite. Since the
regular locus is stable under generalization, we see
that $R'$ is J-1.
\end{proof}

\end{document}
